% arc127 A - Leading 1s の手書きメモの再現（振り返りの文章は thinking.md）
% ビルド: tools/build_thinking.sh arc127/a --preview
\documentclass[a4paper]{article}
\usepackage[margin=1.2cm]{geometry}
\usepackage{handmemo}
\pagestyle{empty}

% 小さい丸囲み（f(i) の表に多い）
\def\hwc(#1,#2){\hwcircle(#1,#2)(0.42,0.42)}

\begin{document}

% ---------------------------------------------------------------- 1 枚目: f(i) の表
\begin{memopage}{20260925\_174906.jpg}
  \hw(1.5,1.5){A --- Leading 1s}
  \hw(1.0,3.0){$x$}
  \hw(3.0,3.0){先頭の1の個数\ \ $f(x)$}
  \redpen(9.6,1.4){$f(x) = \sum_k [\,$先頭 $k$ 桁がすべて 1$\,]$ と\\分ければ、$x$ ごとに数えなくてよい}
  \redarrow(10.2,2.3)(9.4,2.9)

  \hwm(1.0,4.9){\sum_{i=1}^{N} f(i)}
  \hwunread(3.6,4.3){$N$ の上に消した字}

  % i = 1 .. 10
  \hw(0.1,7.4){$i$}
  \hw(0.8,7.4){1} \hw(2.0,7.4){2} \hw(3.4,7.4){3} \hw(5.0,7.4){4} \hw(6.8,7.4){5}
  \hw(9.1,7.6){6} \hw(10.8,7.6){7} \hw(12.3,7.6){8} \hw(13.5,7.6){9}
  \hwstrike(14.7,7.8){10} \hw(16.2,7.8){10}
  \hw(0.0,8.7){\normalsize $f(i)$}
  \hw(1.0,8.7){1} \hwc(1.2,8.7)
  \hw(2.2,8.7){0} \hw(3.3,8.7){0} \hw(4.5,8.7){0} \hw(6.6,9.2){0}
  \hwline(8.7,9.4)(9.2,9.4) \hwline(10.1,9.35)(12.4,9.45) \hwline(12.7,9.5)(13.3,9.5)
  \hw(15.8,9.8){1} \hwc(16.0,9.8)
  \hw(17.2,9.9){2}

  % i = 11 .. 20
  \hw(0.7,10.3){11} \hw(1.8,10.3){12} \hw(3.2,10.3){13} \hw(4.9,10.3){14} \hw(6.5,10.3){15}
  \hw(8.6,10.4){16} \hw(10.6,10.5){17}
  \hwstrike(11.9,10.5){18} \hwstrike(13.2,10.5){19}
  \hw(15.6,10.6){20} \hw(16.8,10.9){+2}
  \hw(0.8,11.5){2} \hwc(1.0,11.5)
  \hw(1.7,11.5){1} \hwc(1.9,11.5)
  \hw(3.05,11.6){1} \hwcircle(3.25,11.6)(0.55,0.55)
  \hw(4.8,11.5){1} \hw(6.6,11.5){1} \hw(8.7,11.6){1} \hw(10.4,11.7){1}
  \hw(13.1,11.9){1} \hwcircle(13.3,11.9)(0.6,0.5)
  \hw(15.8,11.7){0} \hw(17.2,11.5){4}
  \redpen(4.6,12.3){\small $f(11)=2$, $f(12)$〜$f(19)=1$ → 11〜20 の和は 10}

  % i = 21 .. 30
  \hw(0.7,13.3){21} \hw(2.0,13.3){22} \hw(3.1,13.3){23} \hw(4.7,13.3){24} \hw(6.6,13.3){25}
  \hwline(8.6,13.2)(9.0,13.2) \hwline(10.1,13.2)(10.6,13.2) \hwline(10.9,13.2)(11.4,13.2) \hwline(12.2,13.4)(12.5,13.4)
  \hw(15.4,13.5){30} \hw(16.6,13.5){+4}
  \hw(0.7,14.4){0}
  \hwline(1.3,14.4)(1.5,14.4) \hwline(2.2,14.4)(5.2,14.3) \hwline(6.8,14.4)(13.2,14.7)
  \hw(15.3,14.9){0}

  \hw(0.5,15.8){30}
  \hwline(1.3,15.8)(1.6,15.8) \hwline(2.1,15.9)(3.1,15.9)
  \hw(0.6,16.6){\normalsize $\vdots$}

  % 100 .. 110
  \hwunread(0.0,17.5){重ね書き}
  \hw(2.7,17.6){101} \hw(4.0,17.6){102}
  \hwline(5.3,17.5)(6.3,17.4) \hwline(7.6,17.3)(12.6,17.7)
  \hw(12.8,17.8){109} \hw(14.3,17.8){110}
  \hw(0.35,18.6){1} \hwc(0.55,18.6)
  \hw(1.55,18.6){1} \hwc(1.75,18.6)
  \hw(2.85,18.6){1} \hwc(3.05,18.6)
  \hwline(6.1,18.6)(6.4,18.6)
  \hw(12.9,18.9){1} \hwc(13.1,18.9)
  \hw(14.3,18.9){2} \hwc(14.5,18.9)
  \hw(16.4,19.4){6}
  \redpen(4.2,19.6){\small 100〜109 の和は 10、110 まで入れると 12}

  \hw(0.25,19.4){2}
  \hw(0.25,20.1){\normalsize $\vdots$}
  \hw(0.25,20.7){9}
  \hw(15.7,20.9){+9}

  % 1000 .. 1109
  \hw(0.3,21.7){1000,\ 1001} \hw(3.6,21.6){$\cdot$}
  \hwline(4.1,21.5)(8.0,21.5)
  \hw(13.6,21.3){1010} \hw(16.1,21.8){10}
  \hw(0.35,22.7){1} \hwc(0.55,22.7)
  \hw(1.55,22.7){1} \hwc(1.75,22.7)
  \hw(13.7,22.3){1} \hwc(13.9,22.3)
  \hwunread(0.0,23.5){重ね書き}
  \hwline(4.1,23.5)(6.3,23.5) \hwline(7.3,23.6)(11.2,23.6)
  \hw(16.0,23.3){10}
  \hw(16.2,24.1){\normalsize $\vdots$}
  \hw(0.6,24.2){\normalsize 1} \hwcircle(0.75,24.2)(0.3,0.33)
  \hw(1.55,24.2){\normalsize 1} \hwcircle(1.7,24.2)(0.3,0.33)
  \hw(13.2,24.1){1} \hwc(13.4,24.1)
  \hwunread(0.0,25.1){重ね書き}
  \hwline(4.5,25.0)(5.8,24.95) \hwline(7.4,24.95)(8.2,24.95)
  \hw(13.1,25.1){110} \hwunread(14.3,25.7){末尾 8 か 9}
  \hw(15.6,24.9){20}

  \redpen(4.6,22.5){\small 1000〜1009、1010〜1019 は 10、1100〜1109 は 20。\\
    \small 先頭が 1, 10, 11 …の「10 個ずつの塊」を数えている}
  \redarrow(12.4,22.7)(15.8,23.1)
\end{memopage}

% ---------------------------------------------------------------- 2 枚目: 1111〜1300 の数え上げ
\begin{memopage}{20260925\_174909.jpg}
  \hwunread(0.4,0.9){重ね書き}
  \hwunread(1.2,1.8){消した字}
  \hw(4.0,1.0){1111}
  \hwline(5.8,1.2)(6.0,1.2) \hwline(7.0,1.2)(7.7,1.2) \hwline(10.5,1.4)(10.7,1.4)
  \hw(12.7,1.6){1119} \hw(14.8,1.7){120}
  \hw(4.3,1.9){3} \hw(13.2,2.35){3}
  \draw[ink,hand,line width=0.7pt] (3.6,2.3) .. controls (7,1.8) and (11,2.0) .. (14.6,2.1);
  \hwline(6.9,1.8)(7.8,2.3) \hwline(7.0,2.4)(7.9,1.9)
  \hw(15.0,2.8){2} \hwc(15.2,2.8)

  \hw(0.6,3.2){$\cdot$1121\ \ |} \hw(3.7,3.2){1122}
  \hw(9.0,3.35){2+3×8+4 = 6+24=30}
  \hw(14.6,3.9){1130}
  \hw(1.5,4.0){2} \hw(4.0,4.2){2}
  \hwline(5.3,4.3)(6.0,4.3) \hwline(7.1,4.3)(7.8,4.3)
  \hwline(10.4,4.6)(11.5,4.6)
  \hw(9.8,5.1){2×10=20}
  \hwunread(15.0,5.1){2 と記号}
  \redpen(11.5,6.0){\small 1111〜1120 の和（$f(1120)=2$）}
  \redarrow(12.6,5.6)(12.6,3.9)

  \hw(0.5,5.6){1131} \hw(3.4,5.6){\normalsize $\partial$}
  \hwline(4.2,5.8)(11.2,5.9)
  \hw(1.1,6.6){2\ \ -}
  \hwline(4.8,6.8)(9.4,7.1)
  \hw(14.3,6.5){1140} \hw(14.8,7.6){2}
  \hw(11.4,7.8){20} \hw(13.2,8.0){$\cdot$}

  \hw(0.4,8.3){1141}
  \hwline(3.1,8.4)(5.1,8.4) \hwline(5.8,8.3)(6.8,8.3) \hwline(8.4,8.5)(9.3,8.6)
  \hw(1.0,9.3){2}
  \hwline(2.5,9.4)(5.3,9.4)
  \hw(13.8,9.1){1150} \hw(13.9,10.2){2}
  \hw(10.9,10.6){20}
  \hw(0.5,10.4){1151}
  \hw(13.9,11.3){1160}

  \hw(0.4,11.6){1191}
  \hwline(2.4,11.6)(5.0,11.8) \hwline(6.4,12.0)(7.2,12.0)
  \hw(13.6,12.1){1200}
  \hw(0.2,12.4){1201}
  \hwline(2.0,12.5)(6.7,12.9)
  \hw(11.0,12.5){19}
  \hw(13.4,13.0){1300}
  \hw(10.9,13.7){10}

  \redpen(4.0,15.0){1201〜1300 は 100 個とも $f=1$ → 和は 100}
  \redarrow(9.2,14.6)(10.8,14.0)
  \redpen(1.0,17.0){112\_, 113\_, … は 10 個ずつ $f=2$、12\_\_ は 100 個 $f=1$。\\
    塊の大きさは「先頭の 1 の並び」と残りの桁数だけで決まる\\
    → $k$ ごとに「先頭 $k$ 桁が 11…1 の数」を数えれば一般式になる}
  \redpen(1.0,20.4){メモはここ（1300 付近）で止まっている。\\
    桁が増えるたびに区間の切り方が変わり、表からは式にまとまらなかった}
\end{memopage}

% ---------------------------------------------------------------- 図
\clearpage
\section*{図: 先頭 $k$ 桁が 1 の数の区間}
$k$ ごとに、先頭が $k$ 個の 1 で始まる数は $1, 10, 100, \dots$ 個の連続した区間になる。

\begin{center}
\begin{tikzpicture}[x=3.3cm,y=0.7cm]
  % 列 = 桁数 d、行 = k。先頭 k 桁が 1 の d 桁の数は 10^{d-k} 個の区間
  \foreach \d in {1,...,4} {\node at (\d,4.1) {\small \d 桁};}
  \foreach \k in {1,...,4} {\node[anchor=east] at (0.45,4-\k) {\small $k=\k$};}
  \foreach \k/\d/\lab/\n in {1/1/{1}/1, 1/2/{10〜19}/10, 1/3/{100〜199}/100, 1/4/{1000〜1999}/1000,
                             2/2/{11}/1, 2/3/{110〜119}/10, 2/4/{1100〜1199}/100,
                             3/3/{111}/1, 3/4/{1110〜1119}/10,
                             4/4/{1111}/1} {
    \draw[fill=black!8] (\d-0.45,4-\k-0.3) rectangle (\d+0.45,4-\k+0.3);
    \node at (\d,4-\k) {\scriptsize \lab\ (\n 個)};}
  \node[redpen,anchor=west,align=left] at (0.1,-0.9) {\small 各行の塊は $1, 10, 100, \dots$ 個。$x$ が $c$ 行に現れる $\Leftrightarrow f(x) = c$};
\end{tikzpicture}
\end{center}

\end{document}
